\section{Manuscript Overview}
\label{sec:manuscript-overview}

This manuscript is deliberately written in a mathematical register: the encoding and
its metatheory are presented through definitions, typing judgments, and denotations
rather than through the Lean code that realizes them. The artifact behind the thesis
is nonetheless an executable, formally verified program, and a few implementation
concerns are inseparable from its correctness argument. Foremost among these is the
\emph{monadic} structure of the encoder---the way it threads a typing context,
allocates fresh names, and accumulates the emitted SMT script---which we make
explicit where it bears on a design decision or a proof, chiefly in
\Cref{sec:lean-implementation-architecture,sec:binders-names-contexts}, while keeping
the state it manipulates implicit everywhere else.


\paragraph{Notations.}
We fix two notational conventions used throughout the thesis.
We write \leaninl{:=} to denote \emph{definitional assignment}---binding a name to a term,
or giving the body of a definition---as distinct from propositional equality \leaninl{=}
(a proposition asserting that two terms are equal) and from the metatheoretic symbol
\(\triangleq\), read ``defined as'', used in displayed mathematics.

When the proof is unused or irrelevant in a given context, we simply write ``\(\irrelprf\)'' in
such triples, mimicking Lean notations.